
Esta seção atende ao requisito de apresentar um resumo de cada uma das publicações científicas ou técnicas utilizadas.

\section{Publicação 1 -- Heckbert (1986)}

\textbf{Referência:} Heckbert, P. S. \textit{Survey of Texture Mapping}. IEEE Computer Graphics and Applications, 1986.

O trabalho de Paul Heckbert apresenta uma visão geral dos fundamentos de texture mapping. O artigo explica que a técnica permite aumentar a riqueza visual de imagens renderizadas sem exigir que todos os detalhes sejam representados explicitamente por geometria. O autor organiza o problema em torno de dois aspectos principais: o mapeamento geométrico, que relaciona uma textura a uma superfície, e a filtragem, necessária para reduzir artefatos de aliasing.

A publicação é importante para este trabalho porque fornece a fundamentação histórica e conceitual para o uso de coordenadas de textura. O conceito de espaço de textura, representado por \(u\) e \(v\), está diretamente relacionado ao UV mapping utilizado nas seções seguintes. A publicação também ajuda a compreender por que filtragem e amostragem são partes essenciais de uma implementação de textura.

\section{Publicação 2 -- Dong et al. (2020)}

\textbf{Referência:} Dong, J. et al. \textit{Survey of Procedural Methods for Two-Dimensional Texture Generation}. Sensors, 2020.

Dong et al. apresentam uma revisão dos métodos utilizados para geração procedural de texturas bidimensionais. O trabalho parte da importância das texturas para a representação de cenas realistas e discute a geração procedural como um processo baseado em modelos matemáticos.

A publicação é especialmente relevante para o tema deste trabalho porque mostra que texturas podem ser produzidas por procedimentos e funções, em vez de depender exclusivamente de imagens previamente criadas. Isso fornece base para a seção sobre mapeamento procedural e para a formulação \(C=f(P,t)\), utilizada como representação simplificada de uma função que produz o resultado visual a partir de parâmetros espaciais e temporais.

\section{Publicação 3 -- Smelik et al. (2014)}

\textbf{Referência:} Smelik, R. M.; Tutenel, T.; Bidarra, R.; Benes, B. \textit{A Survey on Procedural Modelling for Virtual Worlds}. Computer Graphics Forum, 2014.

Smelik et al. realizam uma revisão de métodos de modelagem procedural aplicados à construção de mundos virtuais. O trabalho aborda diferentes tipos de conteúdo, incluindo terrenos, vegetação, rios, estradas, edifícios e cidades. Entre as vantagens da abordagem estão a possibilidade de gerar grandes quantidades de conteúdo e reduzir a intervenção manual.

Ao mesmo tempo, os autores destacam desafios relacionados ao controle do resultado e à interação do usuário com os métodos procedurais. Essa discussão é importante para este trabalho porque mostra que a proceduralidade não deve ser considerada apenas como uma forma de automatizar a geração: é necessário também permitir controle suficiente para que artistas e desenvolvedores obtenham o resultado desejado.

\section{Relação entre as publicações}

As três publicações complementam-se. Heckbert fornece a base para compreender o texture mapping e os problemas de mapeamento e filtragem. Dong et al. concentram-se na geração procedural de texturas bidimensionais. Smelik et al. ampliam a discussão procedural para a geração de conteúdo em mundos virtuais.

Assim, a literatura utilizada cobre desde o fundamento do mapeamento de texturas até técnicas procedurais e suas aplicações em ambientes gráficos.
